\documentclass[12pt]{article}
\usepackage[utf8]{inputenc}
\usepackage[T1]{fontenc}
\usepackage{hyperref}
\title{Saban\^e Language Resources for K-12 Curriculum Development}
\author{Compiled by Claude Sonnet 5.5 (claude-sonnet-5-5) for Andr\'as Kornai}
\date{}
\begin{document}
\maketitle

\section{Introduction}
Saban\^e (ISO 639-3 sae) is a Nambikwaran language of Brazil. It is critically endangered: as of 2004 there were about three native speakers, all over 50. Only academic descriptions turned up in the search. No school materials were found.

\section{Primary Educational Resources}
None found in the search.

\section{Linguistic and Reference Materials}
\begin{itemize}
\item Gabriel Antunes de Araujo, \emph{A Grammar of Saban\^e, a Nambikwaran Language} (PhD dissertation, Vrije Universiteit Amsterdam, 2004; published by LOT). I found no open-access full text. The SAILS record cites it as the source: \url{https://sails.clld.org/languages/sae}
\item SAPhon phonological inventory for Saban\^e: \url{https://linguistics.berkeley.edu/saphon/en/inv/Sabane.html}
\item Wikipedia overview: \url{https://en.wikipedia.org/wiki/Saban%C3%AA_language}
\item Brazilian journal article that calls for revitalization support (content not checked): \url{https://periodicostestes.bce.unb.br/index.php/ling/article/view/20981}
\end{itemize}

\section{Religious and Community Resources}
None found in the search.

\section{Digital and Audio Resources}
None found in the search.

\section{Curriculum Development Context}
With so few speakers, any curriculum would be a documentation and revitalization effort built on the Antunes de Araujo grammar and in consultation with the community. No evidence of school teaching was found.
\end{document}